
% @author: Sebastian Sepulveda

\documentclass[letterpaper,10.7pt]{article}
\usepackage[spanish]{babel}
\usepackage[utf8]{inputenc}
\usepackage{graphicx}
\DeclareGraphicsExtensions{.jpg,.pdf,.mps,.png}
%Paquetes adicionales, ayudan para portada (algunos)
\usepackage{amssymb}
\usepackage{amsfonts}
\usepackage{amsmath}
\usepackage{fancyhdr}
\usepackage{wrapfig}
\usepackage[dvipsnames]{xcolor}
\colorlet{LightRubineRed}{RubineRed!70!}%https://www.overleaf.com/learn/latex/Using_colours_in_LaTeX
\usepackage{multicol}
\usepackage{changepage}
\usepackage{float}
\usepackage{tcolorbox}
\usepackage{enumitem}
\usepackage{bigstrut}
%https://tex.stackexchange.com/questions/455341/how-to-represent-the-shift-key
%https://tex.stackexchange.com/questions/176398/carriage-return-symbol-new-command
\usepackage{keystroke}
\usepackage{menukeys}
\usepackage{tabularx,ragged2e,booktabs,caption}
\definecolor{gray51}{rgb}{0.51,0.51,0.51}
\definecolor{almond}{rgb}{0.94, 0.87, 0.8}


% Márgenes
\usepackage[vmargin=1.5cm,hmargin=2cm,head=30pt,includeheadfoot]{geometry}

% Interlineado
\linespread{1.5}
\usepackage{hyperref}
\hypersetup{
    colorlinks=true,
    linkcolor=blue,
    filecolor=magenta,
    urlcolor=ForestGreen,
}
\urlstyle{same}
\usepackage{natbib}
\setcitestyle{super}
\usepackage{blindtext}
\linespread{1.0}\selectfont

%%%%%%%%%%%%%%%%%%%%%%%%%%%%%%%%%%%%%%%%%%%%%%%%%%%%%%%%%%%%%
%NUEVOS COMANDOS
%%%%%%%%%%%%%%%%%%%%%%%%%%%%%%%%%%%%%%%%%%%%%%%%%%%%%%%%%%%%%

\newcommand{\hlc}[2][yellow]{ \colorbox{#1}{#2} }

\newcommand{\figref}[1]{\figurename~\ref{#1}}

%%%%%%%%%%%%%%%%%%%%%%%%%%%%%%%%%%%%%%%%%%%%%%%%%%%%%%%%%%%%%
%NUEVOS COMANDOS
%%%%%%%%%%%%%%%%%%%%%%%%%%%%%%%%%%%%%%%%%%%%%%%%%%%%%%%%%%%%%

\begin{document}


\begin{center}
\textbf{{\large Hito 1: Opiniones de otros grupos:}}
\end{center}

\textbf{Total de evaluadores:} 7

\begin{enumerate}
\item \textbf{Motivacion:}\\
Puntaje c/u:  3+2+2+3+2+3+2\\
Puntaje total : 17
\begin{itemize}
	\item Porqué es de interés? En ningún momento lo explican
	\item No hay una introducción
	\item Seria bueno realzar más para qué sería importante el tema
	\item Clara motivacion, hartos ejemplos, falta aclarar porqué esto puede ser interesante
	\item No especifica el porqué es relevante analizar esto, no queda clara su intención
\end{itemize}


\item \textbf{Hipotesis/Problemas iniciales:}\\
Puntaje c/u: 3+2+3+3+3+1+1\\
Puntaje total : 16
\begin{itemize}
	\item No quedo tan claro qué es lo que se desea analizar
	\item Creo que está desordenada
	\item Puede existir cierto sesgo en Amazon con las otras marcas, quizás constractar con otro tipo de e-commerce que abarque distintos tipos de marcas
	\item No queda clara la hipótesis, ésta se infiere, pero no se aprecia textualmente
	\item No explica un problema en concreto, solo se dice el que hará. Esta muy relacionado a la motivacion, si define bien la intención tendrpan una buena hipótesis
\end{itemize}


\item \textbf{Descripcion de los datos que se van a utilizar:}\\
Puntaje c/u: 3+3+3+3+2+2+3\\
Puntaje total : 19\\
\begin{itemize}
	\item Podrian mejorar la visualización de los gráficos
	\item Si un gráfico muestra pocas diferencias entre los valores sería bueno ver otro gráfico que muestre más clara las diferencias de los valores, especialmente si los consideran al analizar
	\item En un gráfico en especial no se podía leer las marcas, usan diferentes gamas de colores por marca, donde algunos colores no ayudan debido a similaridad y el ancho de las barra
	\item ordenar los gráficos, poner titulos descriptivos a los datos, presentar el dataset de manera más entendible
	\item Buen analisis y conclusiones. Cuidado con el scatter plot con regresiones, no se nota o distingue
\end{itemize}


\item \textbf{Calidad general de la presentación:}\\
Puntaje c/u: 2+2+2+3+3+2+3\\
Puntaje total : 17
\begin{itemize}
	\item No todos estaban igual de preparados
	\item Las ideas generales no quedaban muy clara, lo bueno seria dar más énfasis a la idea principal que quieren decir
	\item Mejorar volumen al hablar
	\item Buen manejo del espacio y de los recursos
	\item Falta conexión entre expositores
	\item Dentro del tiempo establecido y se nota preparación
\end{itemize}

\end{enumerate}


\end{document}

\begin{wrapfigure}{r}{0.4\textwidth}
\centering
  \includegraphics[scale=0.45]{images/tarea3/mario_noreparado.png}
  \captionof{figure}{Mario defectuoso}
  \label{fig:mario2}
\end{wrapfigure}


\begin{figure}[!htb]
   \begin{minipage}[h]{0.3\linewidth}
  \centering
  \includegraphics[scale=0.35]{images/tarea3/original_photo.png}
  \captionof{figure}{Image test original}
  \label{fig:colores1}
\end{minipage}
\hfill
   \begin{minipage}[h]{0.3\textwidth}
     \centering
  \includegraphics[scale=0.40]{images/tarea3/or_ph_with_err.png}
  \captionof{figure}{Image test defectuoso (300px)}
  \label{fig:colores2}
   \end{minipage}
 \hfill
 \begin{minipage}[h]{0.3\textwidth}
  \centering
  \includegraphics[scale=0.45]{images/tarea3/with_300_px.png}
  \captionof{figure}{Image test reparada (300px)}
  \label{fig:colores3}
   \end{minipage}
\end{figure}

Al inicio se comenzó obteniendo los vectores de colores de la imagen, lo que nos entrega los valores del dominio sobre el cual vamos a trabajar y las condiciones de borde tipo Dirichlet (colores conocidos del borde del dominio). La forma de obtener estos valores fue recorrer la matriz o paleta de colores hasta encontrar un vector blanco, en ese instante se comienza a consultar a los ``vecinos'' del vector blanco, si son colores distintos a blanco entonces pertenecen a un caso borde y lo guardamos, en caso contrario no hacemos nada. De esta forma obtenemos el \textit{sténcil} de 5 puntos con el cual se puede resolver el problema de Laplace a través de una ecuación matricial de la forma $A \vec{u} = \vec{b} $. Dado a que en las condiciones del problema no se especifica claramente, se asumió que el sector desconocido podría estar en una esquina o en un borde de la imagen original, por lo que el algoritmo hace la consulta de los colores vecinos a cada color desconocido respecto a un sector que pertenezca a la imagen, sin superar los bordes de la matriz de la imagen. Otro problema que se presentó en esta etapa fue el largo que tenía el vector de colores, que podía ser de largo 3 o 4. Por ende, se implementó una condición sobre el largo del vector color (el cual sería una condición de borde) para que se pudiera sumar con otros colores (otras condiciones de borde) sin problema.


Finalmente, con el resultado de la matriz $A$ y conociendo el vector de las condiciones de bordes $b$, podemos resolver el problema con la función de \texttt{numpy:} \texttt{np.linialg.solve} separando los valores RGB de cada elemento en el vector $b$, y luego uniendo el resultado de los tres resultados, reemplazandolo en el sector desconocido de la imagen. Como resultado se entrega la imagen en formato \texttt{.png}\\

La resolución para construir la matriz $A$ que representa los coefientes de las incógnitas es posible realizarla con una matriz de ceros e ir reemplazando por valores númericos según el indice de cada elemento desconocido. Sin embargo, este método se vuelve extremadamente complejo a medida que la imagen va creciendo y que el dominio de las incognitas se vuelve más pequeño, voviendo la cantidad de ceros extremadamente mayor a la cantidad de elementos distintos de ceros en la matriz.

\begin{figure}[!htb]
   \begin{minipage}[h]{0.3\linewidth}
  \centering
  \includegraphics[scale=0.45]{images/tarea3/mario_originalzoom.png}
  \captionof{figure}{Mario original}
  \label{fig:mario1}
\end{minipage}
 \hfill
 \begin{minipage}[h]{0.3\textwidth}
     \centering
  \includegraphics[scale=0.45]{images/tarea3/mario_reparado_zoom.png}
  \captionof{figure}{Mario reparado}
  \label{fig:mario3}
   \end{minipage}
\end{figure}

Esto motiva a tener una estructura de datos para matrices capaz de almacenar solo los valores que no tengan $0's$. Esta estructura es llamada una \textit{sparse matrix}, la cual es utilizada con este fin
00  01  02  03  04 05
10  11 -12  13  14 15
20 -21  22 -23 -24 25
30  31 -32 -33  34 35
40  41  42  43  44 45
50  51  52  53  54 55

agrego 
sumo con el anterior 1 (sa1)
sa2
sa3
sa4
sa5
salgo y divido por total